%\input ../util/bookmacros.tex
%\input ../util/docparams.tex

\chapter{Differential Calculus}

\section{A Motivating Problem for Differential Calculus}
Let me be clear that there is not only one motivating problem for this 
subject. But I think the problem we introduce illuminates 
the parade of definitions you see in standard calculus books dealing with 
differential calculus. You would in the 
course of study of a standard book come across such a problem only 
after you had built up 
a certain amount of mathematical machinery. As mentioned in the preface,  
I prefer to see the problem 
first and see what motivated the machinery. 
It helps me, and I hope it helps you,
get a better perspective of the subject.
\proclaim Problem.
Suppose one is to make a storage box from a flat 
square sheet of metal 12 inches by 12 inches. The box is made
by doing the following:
\beginEnum
\enum{Cut along the dotted lines; that is, cut out the corner pieces.}
\enum{Fold the remaining four ``flaps'' up. }
\endEnum
Now, the resulting ``box'' is a container without a lid.
See the figure below. \epsfigure{eps/metalbox.eps}{Metal Sheet}

Question: What should the ``flap distance'' $x$ be in order to make a
box with maximum volume?

First, it seems clear that the volume is a function of the parameter $x$. 
What is this function? The volume of a box is the area of the base times 
the height. The area of the base is $(12 - 2x)^2$ while the height of the box 
is $x$. So, the volume function (in cubic inches) can be expressed 
as\footnote{\kern 0.5pt \raise 0.5ex \hbox{\dag}}%
{See the appendix ``Functions, Formulas, Graphs, and Notation''.}
$V(x) = x (12 - 2x)^2$. It is not clear from looking 
at this formula what the best choice of $x$ should be. We could try many 
values of $x$ and see which gives the largest value of $V$; 
it's clear from the diagram
that $x$ must be in the range $[0,6]$ (in inches). 
This is a valid approach, and 
though we may not get an exact best answer it could give us a value close 
to the best. For instance,
we could plug in the input values $0,1,2, \ldots, 6$, in the function $V$ 
seeing which gave the biggest answer. But how would we know if we missed any
input value to try? We could increase the \index{sampling} by 
plugging in the input values $0, 1/2, 1, 1 \, 1/2, \, \ldots, \, 6$ into $V$. 
Again, comparing the values of $V$ after plugging in each of these inputs 
would perhaps give a better answer. But to find 
the best solution we would have to plug in all the numbers in the interval 
$[0,6]$ into the function $V$ to find the best solution. The problem here is 
that there are an infinite number of numbers in the interval, $[0,6]$!

One of the difficulties with this problem is one of visualization. We 
can visualize a box and its construction but the function $V$, or for that 
matter any function, is a fairly abstract notion. It is not a number, it 
is a thing 
that takes a number as an input and produces a number as an output. One way to 
make a function more tangible is to draw its {\it graph\/}. When we do 
this we get 
back to things that we can see and touch. By drawing the graph of $V$, 
we make $V$ more concrete, and we can see another way to find an input, $x$,
which maximizes $V$. See the figure below. 
\epsfigure{eps/boxvolume.eps}
{Box volume as a function of flap distance.}

If you place a level\footnote{\kern 0.5pt \raise 0.5ex \hbox{\ddag}}%
{That is, something with a straight edge like a level or a ruler.} against the graph at any point, the level is 
said to be {\it tangent\/} to the graph at that point. The line that passes 
through the level is called a {\it \index{tangent line}\/}.
If the level is placed at the high point of the 
graph, the level lies flat. That is, the \index{tangent line} at that point is flat.
See the figure below. 
\vskip -.2in
\epsfigure{eps/tangents.eps}{Graph of a function and sample tangent lines.}

So, here is a strategy for finding the \index{maximum}: If someone gives us a 
graph of this function, place a level at some point on the graph 
and move it along the graph. Notice that at a bump in the graph -- where 
potentially there is a maximum or minimum  value -- the level is flat. If we 
mark the point or points where this occurs 
you have a set of input candidates for the maximum. You need 
only compare the value of 
the function to maximize at these points and at the end points
of the interval over which you are searching for a maximum. This, in general,
should be a small number of points. See the figures below.
\epsfigure{eps/humpfunc.eps}{Find the maximum of 
$f(x)$ on the interval [0,4]}
\epsfigure{eps/humpfunc2.eps}{Find the maximum of 
$f(x)$ on the interval [0, 4.2]. Although this is the same function 
as the last figure, the change in the interval leads to a different solution.}

The great thing about this method is that if 
we can identify these points we have reduced the comparison to a handful
(hopefully) of points. 

The problem with this method is that it is not much different from the last 
method. To draw a good approximation to the graph of this function over the interval of interest, 
$[0,6]$, you need to \index{sample} the function
at a large number of points on the interval $[0,6]$. The more you sample, the 
better the approximation to the graph. So, the work is not really much different. To make matters 
worse, for some applications the interval could be huge. To draw the graph 
would be a great deal of work.

Suppose we had a way to numerically measure the ``steepness'' 
of the level applied at any point of the graph 
{\it WITHOUT HAVING TO GRAPH IT.\/}
In addition, this ``steepness measure'' had a numeric value 
of $0$ when the level was flat. In effect, we would have a new function, 
which I will call $V'$ such that $V'(x)$ is the ``steepness'' of $V$ at $x$. 
If someone were 
to hand us this function given $V$, we could find an $x$ which \index{maximizes}
$V$ by examining {\it only\/} points in the set, 
$\{\,x \mid  V'(x) = 0 \,\}$ (the places where $V$ is flat), and 
the end points.
To find a best $x$, we would plug in 
each of these inputs into $V$ -- hopefully a small number of inputs -- and 
find a maximizing $x$. Notice again that instead of having to sift through an 
infinite 
number of input values we would only have to check a small handful. 

But all this is predicated on easily finding the magical function $V'$ given 
$V$. It turns out that the ``measure of steepness'' that we will introduce can 
be computed as $V'(x) = 12 (12 - 8x + x^2)$. Now we need only find where the 
steepness is $0$. But this is a quadratic and can easily be solved. 
In this case, $\{ \, x \mid V'(x) = 0\, \} = \{2,6\}$. If we add the end 
points of the 
interval $[0,6]$, then the set of points to check becomes, $\{0,2,6\}$.
Plugging either $0$, or $6$ into $V$ gives $0$. If we plug 
in $2$ into $V$ we get $128$. Therefore, $x = 2$ gives us the 
biggest value for $V$ 
and is therefore the solution to our problem.

Again the magic is that rather than compare an infinite number of input values 
we end up looking at only a handful of \index{candidate input}s 
that have the potential to produce the maximum value of $V$ provided:

\beginEnum
\enum{We are able to easily find the steepness of the graph of the function 
to maximize as a function of the input, $x$. That is, we are able to find a 
steepness function that has a nice formula.}
\enum{We are able to find the places where the steepness function is zero. That 
boils down to setting the steepness function to zero and solving for
inputs. We then need only compare the values of the function to maximize
on these inputs (and the end points of the interval).}
\endEnum

\section{Measure of Steepness}
How do we measure steepness of these tangent lines, or lines in general? 
Road builders have a similar problem when describing the steepness of their 
roads. The standard way to describe road steepness
is to refer to the road's grade. This is the amount 
of ``rise'' of the road divided by its (horizontal) ``run''%
\footnote{\kern 0.5pt \raise 0.5ex \hbox{*}}{This ratio is actually multiplied 
by 100, but this distracts from the main idea.}. 
The measure is 
{\it independent of the particular ``run'' length\/}. 
This can be seen from the
graph below using simple geometry. This ``rise'' over ``run'' 
number is called the 
slope of a straight line. It has the property that it is zero when 
the line is flat.\epsfigure{eps/simtriangs.eps}{The grade gives the same value
regardless of the amount of run.}

So, our definition of steepness of a line is the slope of the line -- see the appendix ``Straight Line Functions''. 
Therefore,
the steepness at a point on the graph of a function $f$ at $x$ is the 
slope of the line tangent to the graph of $f$ at $x$. Let's formalize this.

\proclaim Definition Steepness. The steepness of a function $f$ at $x$ is
the slope of the line tangent to the graph of $f$ at $x$ 
(provided the tangent exists).

Remember, we need to move from the geometric view of \index{steepness}  (a tangent line)
to an algebraic/mechanical one so we can compute efficiently. That is, we need to avoid having
to graph our function in order to solve the maximization problem.
What we have now done, in effect, is represent our {\it geometric\/}
notion of a \index{tangent} line in an {\it algebraic\/} way as
a single number. This is important enough to state formally.

\proclaim Theorem(Slope Representation). A non-vertical \index{straight line} passing
through the origin can be represented by a {\it single\/} number - its \index{slope}.

{\bf proof:\/} From the appendix on straight line functions we see that 
a line with slope $m$ passing through the origin is the graph 
of the {\it linear function\/} $f(x) = m\,x$. Thus $m$ completely characterizes
such a line.

The big question now is: How do we compute this measure of 
steepness for a given function?

\section{Computing the Slope: Approximating Tangents}
So, we know how to measure the steepness of a straight line - it is its 
\index{slope}. Although it's clear what we mean by 
\index{tangent} when we physically put a \index{level} 
against a graph, it is not so clear how to compute it algebraically
in terms of the original function. 

One way to get a handle on this idea of \index{tangent} is to work with approximate 
tangents to a function at a point, $(x, f(x))$ on the graph. 
In the figure below you see that the line which passes through the points 
$(x, f(x))$ and $(x+h, f(x+h))$ is approximately the same as the 
tangent line. The approximation gets better the closer $h$ is to $0$.
\epsfigure{eps/tangentapprox.eps}{Tangent and approximating tangent.}
Therefore, 
if we can find the slope of the \index{approximating tangent}s then we will get an 
approximation of the slope of the tangent line. Computing the slope 
of the \index{approximating tangent line} is easy, we just need a ``rise'' and a ``run''. 
But the two points on the approximate line provide us with this.
The vertical ``rise'' between the two points is $f(x+h) - f(x)$ while the 
horizontal ``run'' is 
$h$ (we use the variable $h$ to indicate that it is a 
{\it horizontal\/} increment). 
Therefore, the slope of the approximating tangent line is%
\footnote{\kern 0.5pt \raise 0.5ex \hbox{\dag}}{Some books will use $\Delta x$ 
rather than $h$ to indicate a small ``run''. The \index{approximating slope}
is then written as
$\frac{f(x + \Delta x) - f(x)}{\Delta x}$.}
$$
\eqalignno{
{f(x+h) - f(x) \over h} && (1) \cr
}
$$
So, if the approximation to the slope gets better and better the closer
$h$ is to $0$, then it must be best when $h$ is $0$. 
However, if we plug $0$ in for $h$
in the equation above we get $0/0$ -- oops! -- $0/0$ is not defined.
What we need to do is track the values of (1) as $h$
gets smaller and smaller
and see where these values ``tend''.
This is the {\it approximate, then take a limit\/} move of the preface --
the same move we used implicitly in chapter 1 when we replaced a slowly
varying factor by a constant; this time we follow the approximation all
the way to its limit.

For instance, given the sequence of numbers 
$1, 1/2, 1/3, \ldots $ , the $10^{\rm th}$ element in the sequence is $1/10$, 
the $103^{\rm rd}$ element in the \index{sequence} is $1/103$. In general, the $n^{\rm th}$ 
element of the sequence is 
$1/n$. 

Question: Where is the \index{sequence} tending? Put another way, what is the 
{\it limit\/} of the sequence?

Answer: It seems clear that although this \index{sequence} never reaches 0 it is heading
there. 

Before continuing, let us introduce a notational convenience.
To make it easier to talk about sequences we will use the 
notation $\{1/n\}_{n=1}^{\infty}$ 
to refer to the sequence $1, 1/2, 1/3, \ldots$; and more generally, 
we use the notation from chapter 1 to refer to a sequence of values: 
$\{x_n\}_{n=1}^{\infty} = \{x_1, x_2, \ldots\}$.

Now, referring to the \index{sequence}, $1, 1/2, 1/3, \ldots$ 
($\{1/n\}_{n=1}^{\infty}$ in our notation) we 
say that its limit is $0$.
This leads to a general operational definition of a \index{limit}.

\proclaim Definition (Intuitive Limit). We say a \index{sequence} 
$\{x_n\}_{n=1}^{\infty}$ has the \index{limit} $a$ if the 
sequence ``tends to'' $a$ as $n$ ``tends to'' $\infty$. We write 
this with the following notation: $\lim\limits_{n \rightarrow \infty} x_n = a$.
We also say that $x_n$ {\it \index{converge}s\/} to $a$.

We leave for the time being what exactly ``tends to'' means. When Calculus 
was originally developed this notion was not rigorously understood. It 
was defined much the same way beauty is defined in many people's minds:
``I can't define it, but I know it when I see it.'' For a precise definition 
of limit see the appendix ``Infinite Sequences and Limits''.

{\bf NOTE:\/} In our definition of a limit, the limit point $a$ is a number. It can't be infinity or negative infinity.
So given the sequence of numbers $1, 2, 3, \ldots$, this sequence does not have a limit as it doesn't tend to any number.

Now, going back to the \index{approximating slope}s we have our first major 
definition.

\proclaim Definition (Slope of Tangent Line). The \index{slope} of the 
\index{tangent line} at a point $(x,f(x))$ 
of the function $f$ exists if the \index{limit}
$$
\lim_{n \rightarrow \infty} {f(x+h_n) - f(x) \over h_n}
$$
exists for {\it every\/} \index{sequence} $\{h_n\}_{n=1}^\infty$ such that $h_n$ goes to 0 with $h_n \ne 0$ for any $n$;
{\it and\/}, every one of these limits is identical.
In this case, the slope of the \index{tangent line} is this {\it unique\/} limiting value. 

That is, if {\it every\/} sequence of 
approximating slopes tends to the same value, then the slope of the tangent line exists and is this value.

{\bf Discrete Connection:\/} We may relate this ``infinitesimal differencing'' to the discrete
differencing in chapter 1. Fix a step size $h \ne 0$ and sample $f$ at the points
$x, x+h, x+2h, \ldots$; that is, define a sequence, named $y$, by $y_k = f(x + k\,h)$ and
a sequence $t$, by $t_k = x + k\,h$.
The approximate tangent slope $(f(x + h) - f(x))/h$ may
be written: $\frac{\Delta y_0}{\Delta t_0}$.
When the step size is $h = 1$, this is {\it exactly\/} the difference operator $\Delta$
of chapter 1 applied to the samples of $f$ (the ``run'', $\Delta t_0$, is just $1$).
The derivative is $\Delta$ run on a shrinking step: difference, divide by the step,
and let the step tend to $0$.

{\bf Note:\/} This \index{limit} may not exist. There are many examples of functions 
where the limit does not exist at a particular point or a collection of points.
Graphically this means the tangent line doesn't exist or is not unique -- how can this be?
We list a few examples of functions with points on their graphs 
where the tangent line does not exist. 
\epsfigure{eps/discount.eps}%
{This function does not have a \index{tangent line} at $(5,f(5))$.
All of the possible approximating slopes tend to $\infty$ -- not a number.}
\epsfigure{eps/oscillating.eps}%
{The function ${\rm Osc}(x) = \cases{x \sin(1/x) & if $x \neq 0$ \cr
0 & if $x = 0$\cr}$
has no tangent at $(0,{\rm Osc}(0))$.
That is, there are many tangents depending on the sequence of approximating tangents we choose -- but no {\it UNIQUE\/} tangent.}
\epsfigure{eps/nonuniquetang.eps}
{No {\it UNIQUE\/} tangent line at $(0, f(0)) = (0,0)$. 
The two possible tangents according to our definition are the lines with slopes 1 and -1 that go through the origin.
We note in the graphic that one may consider other lines as ``tangents'' -- any line that can ``teeter'' on the graph at $x=0$.
This leads to the notion of a sub-differential -- a set of ``valid'' tangents rather than just one. We will not consider this notion
here, but it is interesting to note that the notion of a derivative is not necessarily set in stone.}

\proclaim Definition (Derivative). The \index{slope} of the \index{tangent}
line (when it exists) at a point $(x,f(x))$ 
is called the {\it derivative\/} of $f$ at $x$ and is denoted $f'(x)$. We take $f'$ to mean the function whose 
value at any $x$ is the derivative of $f$ at $x$. The domain of this new function will be, in general, a subset of the domain
of $f$.\footnote{\kern 0.5pt \raise 0.5ex \hbox{\ddag}}{See the appendix ``Functions, Formulas, Graphs and Notation''.}  
We call $f'$ the \index{derivative} of $f$.
We say that we {\it \index{differentiate}\/} $f$ to get $f'$.

We will also write $(f(x))^\prime$ to denote $f'(x)$ when $f$ is represented by an expression, or if it is more readable.
For instance, $(g(x) h(x))^\prime$ means the derivative of the function represented by the expression of $g$ times $h$
evaluated at $x$.

\section{Recap}

Recall that in our example problems we needed a function $f'$ which would
give the steepness of the \index{tangent line} at each point
$(x,f(x))$ (now called the \index{derivative}).
The idea was that $f'$ would be easy to compute from $f$. From the last 
section, it seems things are getting more complicated. In the process of 
trying to find the slopes of the tangent lines, we introduced a fairly 
complicated notion of limit. Further, it seems we are back to having the 
same problem we had with sampling and graphing the function to find a 
tangent slope. In those methods one needed to sample the function at many 
points. In fact, for a completely accurate result all these methods require
one sample $f$ at an infinite number of points to find the slope of $f$ at 
just one point $x$! Worse, since we want 
the slope of the tangent line not just at a given point $x$, but at all points, 
things look grim. 

The good news is that while the limiting process of finding tangent 
slopes (the \index{derivative}) 
is an infinite process, it turns out that for a large number of nice 
functions, $f$, one can write down a formula for $f'$ at {\it any\/} point without doing 
an infinite amount of \index{sampling}. 

\section{Derivatives of Straight Line Functions}
We start with straight line functions first because the tangent line at any point of a straight line function
is the straight line function itself.
Therefore, we don't have to go through any complicated limiting process.

The simplest such straight line functions are the constant functions.
A \index{constant function} is one which has 
only one output for all inputs. Such a function has the following formula:
$$
f(x) = c
$$
What is the \index{slope} of the \index{tangent line} at a given point $x$? The graph
of this function is a flat line whose height is $c$. Its tangent line 
is the line itself. We know a flat line has $0$ as its slope. 
So, what is the value of the derivative function, $f'$ for each $x$? It is 
another constant function which is identically zero. That is,
$f'(x) = 0$. Again, this says that the slope of the tangent line (derivative)
at $(x,f(x))$ is just the slope of the graph of $f$ at $x$, which has slope $0$.

What about finding the derivative function for a general 
straight line function?
See the appendix ``Straight Line Functions'' for more information 
about the formulas for straight line 
functions. In this case, the function has the formula: $f(x) = m\,x + b$.
As in the constant function case the tangent line
to each point $(x,f(x))$ is just the graph of 
$f$. Therefore, the slope of the tangent line (\index{derivative}) is just the \index{slope} 
of the graph of $f$ which is $m$. So, $f'(x) = m$.

In each case we avoided the \index{limiting process} because the graph of the 
functions {\it were already\/} the \index{tangent}s to the function. 
This is not so for more complicated functions.

\section{Computing a Non-Trivial Derivative}
From algebra class, the next functions studied after 
straight line functions are the quadratic functions. 
Let us choose the simplest 
quadratic, $f(x) = x^2$. Unlike straight line functions, the graph of this function is not 
the same as the tangent line at any point.
To compute the \index{derivative} at a point $x$, 
we will have to use 
the \index{approximating tangent line}s to compute the answer. Let 
$\{h_n\}_{n=1}^{\infty}$ be a sequence of numbers which tend to $0$ and have the property that 
$h_n \ne 0$ for all $n$.
Let us compute, for a given $n$, the slope of the $n{\rm th}$ \index{approximating tangent}. 
This is given by $f(x+h_n) - f(x) \over h_n$ which is
$$
(x + h_n)^2 - x^2 \over h_n
$$
Expanding the expression $(x + h_n)^2$ and simplifying yields:
$$
2x h_n + h_n^2 \over h_n
$$
Since $h_n$ is not $0$, we may cancel an $h_n$ from the top and the bottom 
of the last expression.
This yields
$$
2x + h_n
$$
What does this tend to as $n$ tends to infinity?
The limit is $2x$ since $h_n$ goes to zero. Notice, that this is true for {\it every\/} such sequence 
$\{h_n\}_{n=1}^\infty$ that goes to zero and $h_n \ne 0$.
In other words, we've found the \index{derivative} (slope of the tangent line)
at {\it each point\/} $x$. The formula is $f'(x) = 2x$.
Notice, too, that the expression $2x + h_n$ {\it is\/} the discrete answer:
in chapter 1 we found $\Delta n^2 = 2n + 1$, which is the same expression
with a step size of $1$. The limit simply kills the step-size correction
term, leaving $2x$.
{\it Notice that in one calculation not only have we found a formula for the 
derivative at one particular point, but for all points!}

\section{A Function without a Derivative}
As an example of a function that is not \index{differentiable} at a point,  consider the absolute value 
function, $f(x) = |x|$, from the previous figure. The figure demonstrates that there is not a {\it UNIQUE\/} tangent at $x=0$. We now 
show that the derivative does not exist when $x=0$ according to our definition.
Let $h_n = 1/n$ and try to take  the limit of $\frac{f(0 + h_n) - f(0)}{h_n}$.
This is $\lim_{n\rightarrow \infty} (1/n - 0)/(1/n)$, since $h_n > 0$;  which is $\lim_{n\rightarrow \infty} 1 = 1$.
Now let's try a different sequence converging to zero. Let $h_n = -1/n$, then
$\lim_{n\rightarrow \infty} \frac{f(0+ h_n) - f(0)}{h_n} = (1/n - 0)/(-1/n)$, since $h_n < 0$; this is
$\lim_{n\rightarrow\infty} -1 = -1$. Since different \index{sequence}s converging to $0$ provide different
limiting slopes, the derivative does not exist for this $f$ at $0$.

In fact, one can show that if $w$ is any sequence with $w_n > 0$ such that $w_n$ goes to $0$ as $n$ goes to infinity then
$\lim_{n\rightarrow \infty} {f(0+w_n) - f(0) \over w_n} = 1$. And, if $z$ is any sequence with $z_n < 0$ such that $z_n$ goes
to $0$ as $n$ goes to infinity then
$\lim_{n\rightarrow \infty} {f(0+z_n) - f(0) \over z_n} = -1$.
For the derivative to exist at a given point we need the approximating slopes formed from {\it any\/} sequence $h$
 (with $h_n \ne 0$ such that $h_n$ goes to $0$ as $n$  goes to infinity) to have the same limiting value. 
It is not enough to have it true for one sequence or a family of sequences,
but for  {\it all\/} sequences with the requisite properties.

However, we did find the \index{derivative} of $f(x) = x^2$; so,
one function down, how many more to go? Well, there are an infinite 
number of functions and this was one of the easy ones. This may mean that 
given a function which has nice algebraic properties, we can do manipulations 
like we did above and get a formula for the \index{derivative} at all points of 
the function. This is indeed the case for many functions; the algebra just gets messier for more 
complicated functions. Can we do better than to do this limiting process 
again and again?

If we look at our function $f(x) = x^2$ with derivative function $f'(x) = 2x$
and then look at the related function $g(x) = x^2 + c$, where $c$ is some 
constant, 
can we determine the \index{derivative} function, $g'$, from our knowledge of 
$f$ and its derivative $f'$ without going through the \index{limiting process}? 
Yes, the graph of the function $g$ is just 
the graph of the function $f$ shifted up by the amount $c$. That is, $g$ has 
the same shape as $f$ and so the slope of its tangent lines must be the same 
as $f$. We've just argued that $g'(x) = f'(x) = 2x$ without having to go 
through the limiting process. 
What about other functions which are simple modifications of $f$, can we 
find their derivatives so easily? For instance, consider the function 
$h(x) = 2x^2$. Can we relate the derivative $h'$ to $f'$ without 
going through a \index{limiting process}? The answer again is yes.

\section{Differential Calculus: Computational Strategy}
So, we see \index{differential calculus} comes down to finding the \index{derivative}
function, $f'$, from a function $f$. 
Our plan to find the derivative of a function is: 
Compute the \index{derivative} of 
a base set of functions; we only have to do 
this once. Next, find formulas for computing the derivatives of 
functions composed in various ways from these \index{base function}s. The base 
functions and their compositions will be the only functions for which 
nice formulas will be found for their derivatives.
That is to say, we 
are giving up on finding a formula for the derivative of an arbitrary 
function. This is 
not so bad since, hopefully, the pool of functions which we can \index{differentiate}
will be big enough for the applications we consider. If not, then we will 
have to compute the derivative of a function the hard way; that is, use 
the limit process directly.

\section{A Base -- a Family of Functions and their Derivatives}
We already saw how to find the derivative of the function $f(x) = x^2$ using 
the limit of \index{tangent line approximation}s.
We present formulas for derivatives of the family of {\it power\/} functions, 
$f(x) = x^n$ but omit the details.

\proclaim Theorem(Derivative of Power Functions). If $f$ is the 
power function $f(x) = x^n$ with $n$ 
an integer, 
then the derivative, $f'$, is given by
$$
f'(x) = n x ^{n-1}
$$
provided $x \neq 0$ when $n < 0$.

Notice that when $n$ is zero the function is $f(x) = x^0 = 1$ and its
derivative is $f'(x) = 0$. For a proof of this theorem see the applications section
of the appendix on ``Mathematical Induction''.

Compare this rule with the differences we computed in chapter 1:
$\Delta n^2 = 2n + 1$ and $\Delta n^3 = 3n^2 + 3n + 1$.
In each case the leading term follows the same pattern -- $k\,n^{k-1}$ for
the $k^{\rm th}$ power -- while the straggler terms are step-size corrections,
the very terms the limit kills off. The power function $x^k$ plays the same role
on the continuous board that $n^k$ played on the discrete one, and its rule
of differentiation is the discrete rule with the stragglers gone.
(Recall from the aside on falling powers in the chapter on Summation
Calculus that the discrete board has its own exact version of this rule:
$\Delta n^{\underline{k}} = k\, n^{\underline{k-1}}$.)

\exercise{Show this result in the case when $n=3$. That is, derive the 
derivative of the function $f(x) = x^3$ using the 
definition of the derivative.}

\section{Linearity of Differentiation}
A common way of constructing new functions from others is through what is 
called a {\it linear combination\/}. 
Given functions $f_1, f_2, \ldots, f_n$ a new 
function $f$ is a \index{linear combination} of the functions $\{f_i\}_{i=1}^n$ 
with constants $\{c_i\}_{i=1}^n$ if 
$$
f(x) = c_1 f_1(x) + c_2 f_2(x) + \cdots + c_n f_n(x)
$$
As an example consider $f_i(x) = x^{i-1}$, for $i \in \{1,2,3\}$.
Then the function
$$
p(x) = 2 + 3x + 5x^2
$$
is a \index{linear combination} of the functions $f_1(x) = 1$, $f_2(x) = x$, and
$f_3(x) = x^2$. The coefficients of the combination (the $c$'s)
are $2$, $3$, and $5$.

We are able to form a much larger collection of functions from 
a set of \index{base function}s in 
this way. And, it turns out one can easily compute 
the derivative of the new function 
in terms of the derivatives of its component functions.

\proclaim Theorem(Linearity of Differentiation). If the functions $f'_1$ and $f'_2$ exist and 
$f(x) = c_1 f_1(x) + c_2 f_2(x) $, then the derivative, 
$f'$, of $f$ exists and may be computed  as
$$
f'(x) = c_1 f_1'(x) + c_2 f_2'(x) 
$$

Note that this means that $(c f)' = c f'$. This follows by letting $f_2(x) = 0$ in the above theorem.
This result may be extended to an arbitrary linear combination of functions.
\proclaim Corollary(Linearity of Differentiation). If the functions: $f'_1$, $f'_2$, $\ldots$, $f'_n$ exist and 
$f(x) = c_1 f_1(x) + c_2 f_2(x) + \cdots + c_n f_n(x)$, then the derivative, 
$f'$, of $f$ exists and may be computed  as
$$
f'(x) = c_1 f_1'(x) + c_2 f_2'(x) + \cdots + c_n f_n'(x)
$$

This follows from the previous theorem and mathematical induction.\footnote{\kern 0.5pt \raise 0.5ex \hbox{\dag}}{See the appendix ``Mathematical Induction''.}
The theorem and corollary say that the operation of differentiation is linear. This is directly analogous to the linearity
result for the \index{difference operation}, $\Delta$, of chapter 1. In effect, differentiation (like differencing)
passes through constants and sums.

We show this for the case when $f$ is the sum of two functions; that is $f(x) = c_1 \, f_1(x) + c_2\,  f_2(x)$.
Let $h$ be a sequence with the property that $h_n \ne 0$ and $h_n$ tends to $0$ as $n$ tends to infinity. 
This sequence allows us to form approximating tangent lines whose slopes at $(x, f(x))$ are
$f(x + h_n) - f(x) \over h_n$. This is related to approximating tangent slopes of $f_1$ and $f_2$ as follows:
$$
\eqalign{
{f(x + h_n) - f(x) \over h_n} & = {\bigl(c_1 \, f_1(x + h_n) + c_2\, f_2(x + h_n) \bigr)  - \bigl( c_1\, f_1(x) + c_2 \,f_2(x)\bigr) \over h_n} \cr
& =  c_1 \biggl( {f_1(x + h_n)  - f_1(x)   \over h_n } \biggr) + c_2 \biggl( {f_2(x + h_n) -  f_2(x)\over h_n} \biggr) \cr
}
$$
Taking the limit of both sides of this equation as $h_n$ goes to $0$ gives\footnote{\kern 0.5pt \raise 0.5ex \hbox{\ddag}}%
{Passing the limits through constant multipliers is discussed in the section 
``Computational Strategy for Limits'' in the appendix ``Infinite Sequences and Limits''.}\hfil\break
$\lim_{n\rightarrow \infty} \biggl( { f(x + h_n) - f(x) \over h_n} \biggr)  = c_1 \, f_1'(x) + c_2 \, f_2'(x)$. Since this
limit exists and is the same for {\it all\/} such sequences $h$ the derivative of $f$ at $x$ exists and is this value. That is,
$f'(x) = c_1 \, f_1'(x) + c_2 \, f_2'(x)$.

\example{Let's use this theorem to find the derivative of 
the polynomial $p(x) = 2 + 3x + 5 x^2$.}

The theorem says that if we 
can {\it recognize\/} $p$ as a \index{linear combination} of functions for which we
know how to compute the derivative, then we can easily compute $p'$.
We see that 
$p(x) = 2 f_1(x) + 3 f_2(x) + 5 f_3(x)$, with $f_1(x) = 1$, 
$f_2(x) = x$, and $f_3(x) = x^2$. Therefore, 
$p'(x) = 2 f'_1(x) + 3 f'_2(x) + 5 f'_3(x) = 2 * 0 + 3 * 1 + 5 * 2x = 3 + 10x$.

The linearity theorem combined with the rule for \index{differentiating powers}
allows us to compute the derivative of any polynomial!

\proclaim Theorem(Polynomial Differentiation). Let 
$p(x) = c_0 + c_1 x + c_2 x^2 + \cdots + c_n x^n$, ($n$ a non-negative integer) then 
$$
p'(x) = c_1 + 2 c_2 x + 3 c_3 x^2 + \cdots + n c_n x^{n-1}
$$

This result follows easily because of the \index{linearity of differentiation}; we need only know how 
to find the derivative of each term. But each term is just a power function whose derivative we know.

Armed with our ability to compute derivatives for our \index{base family} -- the \index{polynomial}s, let's 
go back and solve the original problem we posed about the metal box container. 


\section{Application Revisited}
Consider again the original problem of the metal box container. The 
problem was to maximize the volume of a box made from a piece of sheet 
metal. The volume of the box was a function of 
the ``flap distance'' $x$. A formula for the volume was 
$V(x) = x (12 - 2x)^2$.

We will solve this again using our knowledge of how to compute derivatives. 
We proceed by finding all solutions to $V'(x) = 0$ combined with the end 
points of our search 
and then determine a value which maximizes $V$. 

Now, $V$ is a \index{polynomial} since,
after expanding the formula for $V$ we have: 
$V(x) = x(144 - 48x + 4x^2) = 144x - 48x^2 + 4x^3$.
But we know now how to compute $V'$ when $V$ is a polynomial. According 
to the last theorem it is
$V'(x) = 144 - 2 * 48 x + 3 * 4 x^2$; or,
$V'(x) = 144 - 96 x + 12 x^2$.
Recall the point of having $V'$: {\it By finding those inputs, $x$, such that 
$V'(x) = 0$ we have a list of places where $V$ is flat and therefore a 
potential maximum.\/}
Solving the equation, $V'(x) = 0$, for $x$ implies that $x$ satisfies
$144 - 96 x + 12 x^2 = 0$. Or, $12 - 8 x + x^2 = 0$. Again, the two solutions
are $2$ and $6$. To these values we merge the set of end points, $\{0,6\}$, 
to give the candidate input set, $\{0,2,6\}$. Plugging in $0$ and $6$ 
yields $0$. 
Plugging in $2$ produces $128$. Therefore, $x=2$ is the input that gives 
the maximum value of $V$, the same result as before.

One may argue that with the advent of computers it might be easier 
to find this solution by brute force. This is a reasonable argument.
However, what if someone asked to maximize the volume that could be produced 
from square sheets of metal other than 12 inches? One could resolve 
this problem using Calculus with a parameter representing the 
size of the square. The resulting formula would allow you to solve
an entire {\it family\/} of problems. This is something the brute
force method cannot do.

\exercise{A farmer has $100$ feet of fencing and wants to enclose a rectangular
pen against a straight river bank; the side along the river needs no fence.
Express the enclosed area as a function of the side length perpendicular to the
river, and use the candidate-input strategy of this chapter to find the
dimensions giving the largest area.}

\section{Enlarging the Class of Differentiable Formulas}
For our derivative technique to be useful for more complicated maximization 
problems we need to enlarge the class of functions, $f$, for which we can 
find a formula for $f'$.
As we know from our strategy we need to do two things: Enlarge the set 
of \index{base function}s for which we know how to compute the derivative;
and, find ways to combine these functions in a way that we can determine the 
derivative from the component functions. 
How else can we form more complex functions, 
$f$, from simpler functions, $g$, and $h$? Here are three ways:

\beginEnum
\enum{$f(x) = g(x) h(x)$}
\enum{$f(x) = g(x) / h(x)$}
\enum{$f(x) = g(h(x))$}
\endEnum

It turns out one can find the derivative of functions formed in the ways
enumerated above in terms of the derivatives of the functions used to compose
them. This is just what we need: a function built up from simple pieces --
like $(1 + x^2)^7$, which we will meet shortly -- would be hopeless to
differentiate from the definition alone; the rules below let us differentiate
combinations in terms of their parts.
Here are the three corresponding theorems for computing their derivatives.
We also restate the theorem on \index{linear combination}s for completeness.

\proclaim Theorem(\index{Linearity Rule}). If $g_1'(x), g_2'(x), \ldots, g_n'(x)$ exist and
$f(x) = c_1 g_1(x) + c_2 g_2(x) + \cdots + c_n g_n(x)$, then $f'(x)$ exists and
$$
f'(x) = c_1 g_1'(x) + c_2 g_2'(x) + \cdots + c_n g_n'(x)
$$

\proclaim Theorem(\index{Product Rule}). If $g'(x)$ and $h'(x)$ exist and
$f(x) = g(x) h(x)$, then $f'(x)$ exists and
$$
f'(x) = g'(x) h(x) + g(x) h'(x)
$$

\proclaim Theorem(\index{Quotient Rule}). If $g'(x)$ and $h'(x)$ exist,
$h(x) \neq 0$, and $f(x) = g(x) / h(x)$, then $f'(x)$ exists and
$$
f'(x) = {g'(x) h(x) - g(x) h'(x) \over h^2(x) }
$$

\proclaim Theorem(\index{Chain Rule}). If $h'(x)$ and $g'(h(x))$ exist and
$f(x) = g(h(x))$, then $f'(x)$ exists and
$$
f'(x) = g'(h(x)) h'(x)
$$

To make the chain rule plausible, write $u = h(x)$ and $y = g(u) = f(x)$ and
think of an approximating slope as a ratio of small changes; then
$\frac{\Delta y}{\Delta x} = \frac{\Delta y}{\Delta u} \frac{\Delta u}{\Delta x}$,
and letting the step shrink, the two factors on the right tend to $g'(h(x))$
and $h'(x)$. A careful proof, however, needs the limit machinery of the
appendix ``Infinite Sequences and Limits''.

We derive the product rule as it is similar to the product rule in the Discrete Calculus.
Let $h$ be a sequence with $h_n \ne 0$ such that $h_n$ goes to $0$ as $n$ goes to infinity. 
Given that $f(x) = g(x) h(x)$, consider the slope of an approximating tangent to $f$ at $x$ based on the sequence $h$: ${f(x+h_n) - f(x) \over h_n}$.
There is a naming problem here. We have been using a sequence named $h$ in our approximating slopes but the function $f$ 
is composed of a function with the name $h$. What do we do? Well, we could press on and not let it bother us 
because a subscripted $h$ is an element of the {\it sequence\/} $h$
while $h(x)$ means the {\it function\/} $h$ evaluated at $x$. Or, to be more clear we could change the name of either the sequence or
the function $h$.
We will change the name of the sequence.
There is nothing special
in the name of this sequence. All that is important is that its elements are not $0$ and that the sequence tends to $0$. 
So, let's rename the sequence $h$ to $z$.
An approximating tangent slope of $f$ at $x$ is $f(x+z_n) - f(x) \over z_n$. We now try to relate
this to the approximating tangent slopes of $g$ and $h$.\footnote{\kern 0.5pt \raise 0.5ex \hbox{\dag}}%
{The term $g(x)h(x+z_n)$ is subtracted and added in an
analogous way to the derivation of the product rule in the Discrete Calculus.}
$$
\eqalign{
  {f(x+z_n) - f(x) \over z_n} & = {g(x+z_n) h(x+z_n) - g(x) h(x) \over z_n} \cr
& = { g(x+z_n) h(x+z_n) - g(x) h(x+z_n) + g(x)h(x+z_n) - g(x) h(x) \over z_n} \cr
& = \biggl( { g(x+z_n) - g(x) \over z_n } \biggr)  h(x+z_n) + g(x) \biggl( { h(x+z_n) - h(x)  \over z_n} \biggr) \cr
 }
$$
Taking the limit as $n$ tends to infinity of the equation above gives:\hfil\break
$\lim_{n\rightarrow \infty} \biggl( { f(x+z_n) - f(x) \over z_n} \biggr)  = g'(x) \lim_{n\rightarrow \infty} h(x+z_n) + g(x) h'(x)$.
We are assuming that the limit of $h(x+z_n)$ exists and are using the properties of the limiting operation as described in
the section ``Computational Strategy for Limits'' in the  appendix, ``Infinite Sequences and Limits''. It turns out that a
differentiable function is a continuous function and this means\footnote{\kern 0.5pt \raise 0.5ex \hbox{\ddag}}%
{Also discussed in the section ``Computational Strategy for Limits'' in the appendix ``Infinite Sequences and  Limits''.}
that $\lim_{n \rightarrow \infty} h(x+z_n) = h(\lim_{n\rightarrow \infty} (x + z_n)) = h(x)$. Finally, since this limit exists
for {\it all\/} such sequences $z$, the derivative of $f$ exists and is equal to this value. That is,
$f'(x) = g'(x) h(x) + g(x) h'(x)$.

So, given a function $f$, we can find its derivative if we can
recognize it as being either:

\beginEnum
\enum{One of the base functions for which we know the derivative.
For us, this is the set of all \index{polynomial}s and power functions.}
\enum{Composed (in one of the ways above) of functions for which
one knows the derivatives.}
\endEnum
Keep in mind, there may be many ways to go about computing derivatives. 
It's really a kind of \index{pattern recognition}.

\example{Suppose $f(x) = (1 + x + 5x^2) (1 - x^2)$. Find $f'(x)$.}

We first check to see if this function is one of our base functions.
Currently, we only recognize the set of \index{polynomial}s as our base set.
So, this is not
a base function. Next, we check to see if $f$ can be decomposed into one of
the above patterns. Yes, $f$ can be written as the product of
two functions, $g$ and $h$, where
$g(x) = (1 + x + 5x^2)$ and $h(x) = 1 - x^2$. Do we know how to compute the 
derivative of $g$ or $h$? Yes, each of these are polynomials and we have 
a formula for computing their derivatives. So, using the \index{product rule} we have
$$
\eqalign{
f'(x) & =      g'(x)         h(x) +        g(x)         h'(x) \cr
      & = (1 + x + 5x^2)' (1 - x^2) + (1 + x + 5x^2) (1 - x^2)' \cr
      & = (0 + 1 + 2 * 5x) (1 - x^2) + (1 + x + 5x^2) (0 - 2x) \cr
      & = (1 + 10x) (1 - x^2) - 2x (1 + x + 5x^2) \cr
}
$$
Of course, another way to compute the derivative would be to multiply out 
the product of $g$ and $h$. This would give us a polynomial and we could 
use the formula for differentiating polynomials.

\example{Let $f(x) = 1 + x + x^2 - 2x$. Find $f'(x)$.}

We can find 
the derivative in one of two ways. First, $f$ is a \index{polynomial} isn't it? 
Well, no, not in the form our theorem on polynomial differentiation requires. Our theorems refer 
to polynomials as being a \index{linear combination} of {\it DISTINCT\/} power functions.
In the case of $f$, the ``x'' 
term appears twice; first as $x$ and then later as $-2x$. It is easy enough 
to turn this into a ``canonical'' \index{polynomial}, just add like powers. 
$f$ becomes: $f(x) = 1 - x + x^2$. Now we can apply the theorem for 
computing derivatives of polynomials giving us: $f'(x) = -1 + 2x$.
Another way to compute this is to use the linearity theorem which just 
distributes the derivative through sums and products with constants. It says
$f'(x) = 1^\prime + x^\prime + (x^2)^\prime - 2x^\prime$. 
Applying the rule for differentiating powers and constants 
gives us $f'(x) = 0 + 1 + 2x - 2 = -1 + 2x$. This agrees with our 
previous computation.

These rules are usually easy to apply. The only one which is harder to recognize
is the chain rule. Here is an example of how it is used. Let 
$f(x) = (1 + x^2)^7$. What is the derivative of $f$? Again, in this case we 
could raise $(1 + x^2)$ to the $7^{\rm th}$ power to give us a 
polynomial and then
write down its derivative. We can also recognize $f$ 
as a composition of the function $g(x) = x^7$ with the function 
$h(x) = 1 + x^2$. That is, $f(x) = g(h(x))$. Applying the chain rule for the 
composition of two functions gives
$$
\eqalign{
f'(x) & = g'(h(x))    h'(x) \cr
      & = 7(1+x^2)^6 (1+x^2)' \cr
      & = 7(1+x^2)^6(0 + 2x) \cr
      & = 7 (1 + x^2)^6 2x  \cr
      & = 14x (1 + x^2)^6 \cr
}
$$
Notice that $g'(x) = 7x^6$ and therefore, $g'(h(x)) = 7h(x)^6$. Since
$h(x) = (1 + x^2)$ it follows that $g'(h(x)) = 7 (1 + x^2)^6$.

\exercise{Find $f'(x)$ for each of the following: (1) $f(x) = (1 + x^3)^{10}$;
(2) $f(x) = x^2 \, (1 + x)^5$; (3) $f(x) = {x \over 1 + x^2}$. In each case,
first decide {\it which\/} rule applies.}

\section{Enlarging the Base}
So far, our collection of functions for which we know the derivative is the collection of power functions, $f(x) = x^n$, for integer
$n$, and all polynomials.
Why enlarge it? Because nature keeps handing us functions that are not
polynomials. Growth processes -- money compounding in a bank account,
a population of bacteria doubling again and again -- force exponential
functions on us, and the chapter ahead on differential equations will
need their derivatives.
We now enlarge this base collection to include formulas for computing the
derivatives of \index{logarithmic function}s (and their inverses) as well as 
\index{trigonometric function}s. The next subsections help to explain the number $e$ that 
you may have heard of as well as the angle measure called \index{radian}s.
This section is not really needed to understand the essential aspects
of Calculus; you may safely skip the next two subsections if these functions don't interest you.
You will, however, need to know about the log and exponential functions for the chapter on differential equations.

\subsection{Derivatives of Log and Exponential Functions}
Let's try to compute the derivative at a given point $x$ of the log 
function base $b$, $\log_b$. The logarithm only makes sense for 
positive inputs, so 
throughout our computations, $x$ is assumed positive.
The derivative of this function at $x$ is the limit --
if it exists -- of $\lim_{n \rightarrow \infty} \bigl(\log_b(x+h_n) - \log_b(x)\bigr) / h_n$.
Here $\{h_n\}_{n=1}^\infty$
is a \index{sequence} which converges to $0$ with the property that $h_n \ne 0$ for all $n$.
This may be written as $\bigl(\log_b\bigl(x (1 + h_n/x)\bigr) - \log_b(x)\bigr) 
/ h_n$. Using 
the property of logs, this may be written: 
$\bigl( \log_b(x) + \log_b(1 + h_n/x) - \log_b(x) \bigr) / h_n$; or,
$\log_b(1 + h_n/x) / h_n$. Using the property of logs one more time we may write 
this as $\log_b(1 + h_n/x)^{1/h_n}$; or, 
$\log_b\left[(1 + h_n/x)^{x/h_n}\right]^{1/x}$. We need to find the limit of this  as
$n$ goes to $\infty$.
Since
the logarithm and the power function $z(s) = s^{1/x}$ are continuous for $x > 0$, we may pass the 
limit through these functions\footnote{\kern 0.5pt \raise 0.5ex \hbox{\ddag}}%
{A continuous function, $f$,
has the following property: If $\lim\limits_{n\rightarrow \infty} x_n = x$ then
$\lim\limits_{n\rightarrow \infty} f(x_n) = f(\lim\limits_{n\rightarrow \infty} x_n)
= f(x)$. We say in this case that the limit passes through $f$. See again, appendix C 
for a discussion of continuous functions.} yielding 
$\log_b\left[\lim_{n\rightarrow \infty}(1 + h_n/x)^{x/h_n}\right]^{1/x}$. It turns out that
the limit
inside the \index{log} (regardless of $x>0$)  is the same as the limit:
$\lim_{n \rightarrow \infty}(1 + 1/n)^n$ -- a step we take on faith here;
its justification belongs to the theory of limits (see the appendix
``Infinite Sequences and Limits''). This limit exists and its limiting
value
is called $e$ which is approximately $2.718281828459045$. So the limit becomes
$\log_b\bigl(e^{1/x} \bigr)= (1/x)\log_b(e)$. Therefore, we have the formula:
$$
\bigl(\log_b(x) \bigr)^{\prime}= (1/x)\log_b(e)
$$
Notice the messy constant $\log_b(e)$ tagging along; it depends on our choice
of base $b$. The number $e$ is precisely the base for which this constant
equals $1$, so that the derivative comes out clean -- we will make the same
kind of choice later when we pick radians as our angle measure.
If we choose the base of the \index{logarithm} to be $e$, then
$$
\bigl(\log_e(x)\bigr)^{\prime} = (1/x)\log_e(e) = 1/x
$$
The function, $\ln$, is the name used to denote $\log_e$. 

We can use this knowledge to compute the derivative of $\ln$'s inverse, 
the function $e^x$.
By definition, $\ln(e^x) = x$, taking the derivative of this equation yields 
$\bigl(\ln(e^x)\bigr)' = x' = 1$.
Or,
$$
1 = \bigl(\ln(e^x)\bigr)^\prime = \obrace{\frac{1}{e^x}}{$h'(g(x))$}\obrace{ \vphantom{\frac{1}{e^x}} (e^x)^\prime}{$g'(x)$} 
$$
The last equality follows from the chain rule with $h(x) = \ln(x)$ and $g(x) = e^x$.
Rewriting, this becomes:
$$
(e^x)^\prime = e^x
$$
We repeat this procedure for $b^x$, where $b$ is any positive number. Taking the 
derivative of the identity $\log_b(b^x) = x$, gives
$\bigl(\log_b(b^x)\bigr)' = x' = 1$.
Or,
$$
1 = \bigl(\log_b(b^x)\bigr)^\prime = \obrace{\frac{1}{b^x}\log_b(e)}{$h'(g(x))$} \obrace{ \vphantom{\frac{1}{b^x} }(b^x)^\prime}{$g'(x)$}
$$
The last equality follows from the chain rule with $h(x) = \log_b(x)$ and $g(x) = b^x$.
Rewriting, this becomes: $(b^x)^\prime = \frac{1}{\log_b(e)}b^x$.
Since, $1/\log_b(e) = \log_e(b) = \ln(b)$, we may write this last expression as
$$
(b^x)^\prime = \ln(b) b^x
$$
We finish this section with a formula for computing the derivative of the function $e^{a\, x}$.
This is really the composition of functions $h(x) = e^x$ and $g(x) = a\, x$; that is, 
$e^{a\, x} = h(g(x))$. By the chain rule the derivative is $h'(g(x))\,g'(x)$. Since 
$h'(x) = (e^x)' = e^x$ and $g'(x) = a$ it follows that $(e^{a\, x})^{\prime} = e^{a\, x} a = a \, e^{a \, x}$.
This formula will be of use when we get to the chapter on differential equations.
We list our main results:\footnote{\kern 0.5pt \raise 0.5ex \hbox{\dag}}{The domain of $\ln'$ is the same as $\ln$ which is $\{ x \,|\, x > 0\}$.}
\beginEnum
\enum{$\ln'(x)  = {1 \over x}$}
\enum{$\bigl(e^x\bigr)^\prime = e^x$}
\enum{$\bigl(e^{a\, x}\bigr)^\prime = a\, e^{a\, x}$}
\endEnum
\vskip .2in

\subsection{Derivatives of Trig Functions}
Let us attempt to compute the derivative of the sine function, $\sin$, 
at a point $x$. Again, let $\{h_n\}_{n=1}^\infty$ be a sequence which converges to $0$ with the
property that $h_n \ne 0$ for any $n$. For any $n$, the approximating slope to the tangent line is
$$
\eqalign{
{\sin(x+h_n) - \sin(x) \over h_n} & = {\obrace{\sin(x)\cos(h_n) + \cos(x)\sin(h_n)}{$\sin(a+b) = \sin(a)\cos(b) +
    \cos(a)\sin(b)$}  \mskip -30mu - \sin(x) \over   h_n}  \cr
& = \sin(x) {\cos(h_n) - 1 \over h_n} + \cos(x){\sin(h_n) \over h_n}
}
$$
We need to be able to find the limits (if they exist) of $\frac{\sin(h_n)}{h_n}$ and 
$\frac{\cos(h_n) - 1}{h_n}$. We will have to do some real work to compute the limit 
of $\frac{\sin(h_n)}{h_n}$; but once we do, the other limit, $\frac{\cos(h_n) - 1}{h_n}$ 
will follow and, therefore, we will be able to compute the derivative of the 
sine function. Before we make a precise argument for the limit of 
$\frac{\sin(h_n)}{h_n}$, note in the figure below that when $\theta$ is small,
the length of the segment BC and the length of the arc DC seem to be very close to the 
same value. In fact, one might conclude that in the limit their ratio 
$\frac{\sin(\theta)}{\theta}$ is $1$. This is what we show now.

Examining the figure and comparing the areas of the triangular region ABC to 
the sector ADC to the triangle ADE we find that 
$$
\obrace{\frac{\cos(\theta)\sin(\theta)}{2}}{Area of triangle ABC} \le 
\obrace{\frac{\theta}{2}}{Area of sector ADC} \le 
\obrace{\frac{1 * \tan(\theta)}{2}}{Area of triangle ADE}
$$
\epsfigure{eps/sinderiv.eps}{Comparison of triangle ABC, sector ADC, 
and triangle ADE.}
Here, we are assuming $\theta$ is positive, measured in radians, and 
is less than $\pi/2$. The area of each triangle is found using the formula, 
base * height / 2, while the area of the sector is determined by the formula 
$\theta r^2/2$. In our case, $r = 1$, giving $\theta/2$; again, 
provided $\theta$ is measured in \index{radian}s. We briefly describe 
radian measure for those who have not seen it before.

The radian measure of an angle is the arc 
length of the angular sector 
formed by this angle on the unit circle. That is, if a sector has angle $\theta$ in 
radians, then the corresponding arc (on the unit circle) has length $\theta$ (see 
the figure above).
The arc length of a complete circle 
(360 degrees) is $2\pi$; therefore, there are $2\pi$ radians in 360 degrees. This means
that there are ${360 \over 2 \pi }= {180 \over \pi}$ (which  is approximately 57) degrees per radian. 
Also, the area of the unit circle is $\pi$; so, 
the area of a sector of $\theta$ radians (on the unit circle) is a fraction of the area of the 
full circle $\pi$. What fraction? Well, the sector of $\theta$ radians has an arc length of $\theta$ while
the arc length of a complete circle is $2 \pi$. Since the area of a complete circle is $\pi$, the area
of the sector must be ${\theta \over 2 \pi}$ of $\pi$. So, the area of the sector is $\theta / 2$.
One can use other angular measures (such as degrees or grads) 
but the formulas for the area of a sector have messy constants; and, in turn, 
the derivatives of the trig functions using these angular measures also 
have messy constants. Note that this is not unlike the derivative 
of the logarithmic functions; any choice of base other than $e$ leads to a messier 
formula.

Continuing with our calculations, the first inequality gives 
the relation:
$$
\frac{\sin(\theta)}{\theta} \le \frac{1}{\cos(\theta)}
$$
while the second inequality yields the relation:
$$
\cos(\theta) \le \frac{\sin(\theta)}{\theta}
$$
Combining these gives:
$$
\cos(\theta) \le \frac{\sin(\theta)}{\theta} \le \frac{1}{\cos(\theta)}
$$
A similar argument shows that the relation holds when $\theta$ is negative and larger 
than $-\pi/2$. We may replace $\theta$ with the $n^{\rm th}$ element of our sequence $h$ giving:
$$
\cos(h_n) \le \frac{\sin(h_n)}{h_n} \le \frac{1}{\cos(h_n)}
$$
Now, since the limit of the first and third terms exists and is $1$, it is the case that the middle term has 
the same limit.\footnote{\kern 0.5pt \raise 0.5ex \hbox{\dag}}{This follows from the Squeeze Limit Theorem from the appendix
  ``Infinite Sequences and Limits''.} Therefore,
$$
\lim_{n\rightarrow \infty} \frac{\sin(h_n)}{h_n} = 1
$$
We can use this result to compute the other limit of interest, 
$(\cos(h_n) - 1)/h_n$. First note that 
$$
\frac{\cos(h_n) - 1}{h_n} = \frac{(\cos(h_n) - 1) (\cos(h_n) + 1)}{h_n (\cos(h_n) + 1)}
= \frac{\cos^2(h_n) - 1}{h_n(\cos(h_n) + 1)}
$$
This may be written 
$$
\frac{-\sin^2(h_n)}{h_n(\cos(h_n) + 1)} = -\frac{\sin^2(h_n)}{h_n^2} \frac{h_n}{\cos(h_n) + 1}
$$
Since the last expression is a product of expressions which have limits, we may 
write:\footnote{\kern 0.5pt \raise 0.5ex \hbox{\ddag}}{This follows from the Product of Limits Theorem in the appendix ``Infinite Sequences and Limits''.}
$$
\lim_{n\rightarrow \infty}\frac{\cos(h_n) - 1}{h_n} = -\lim_{n\rightarrow \infty}
\frac{\sin^2(h_n)}{h_n^2} \lim_{n\rightarrow \infty}\frac{h_n}{\cos(h_n) + 1}
= -(1^2)
 * (0 / 2) = -1 * 0 = 0
$$
Applying these limits to calculate the derivative of sine we have:
$$
\eqalign{
\lim_{n \rightarrow \infty} \bigl(\sin(x+h_n) - \sin(x)\bigr) / h_n & = 
\sin(x) \lim_{n\rightarrow \infty}\frac{(\cos(h_n) - 1)}{h_n} + 
\cos(x) \lim_{n \rightarrow \infty} \frac{\sin(h_n)}{h_n} \cr
& = \sin(x) * 0 + \cos(x) * 1 \cr
& = \cos(x)
}
$$
Therefore, we have shown that $\sin'(x) = \cos(x)$. 

The derivative of the cosine function can be done now without much effort.
$$
\eqalign{
{\cos(x+h_n) - \cos(x) \over h_n} &=
{\obrace{\cos(x)\cos(h_n) - \sin(x)\sin(h_n)}{$\cos(a+b) = \cos(a)\cos(b) - \sin(a)\sin(b)$} \mskip -30mu - \cos(x) \over
  h_n }\cr
&= \cos(x){\cos(h_n) - 1 \over h_n} - \sin(x){\sin(h_n)  \over h_n} \cr
}
$$
Taking the limit we 
have $\cos'(x) = \cos(x) * 0 - \sin(x) * 1 = -\sin(x)$.

The other trig functions can be described in terms of $\sin$ and $\cos$ so 
we may find their derivatives through the derivatives of $\sin$ and $\cos$.
We list the derivatives of these trig functions and restate the results for $\sin$ and $\cos$:
\beginEnum
\enum{$\sin'(x) = \cos(x)$}
\enum{$\cos'(x) = -\sin(x)$}
\enum{$\tan'(x) = \sec^2(x)$}
\enum{$\sec'(x) = \sec(x) \tan(x)$}
\enum{$\csc'(x) = -\csc(x) \cot(x)$}
\enum{$\cot'(x) = -\csc^2(x)$}
\endEnum
To see how these formulas are derived, we compute the derivative of the 
tangent function, $\tan$. Since $\tan$ is the ratio of $\sin$ to $\cos$ we can 
use the quotient rule to compute its derivative. The \index{quotient rule} says
$(f/g)' = (f'g - f\,g')/g^2$. In our case, $f = \sin$ and $g = \cos$. Plugging 
these into the \index{quotient rule} gives
$$
\eqalign{
\tan'(x) &= \bigl(\sin(x)/\cos(x)\bigr)' \cr
&= \bigl(\sin'(x) \cos(x) - \sin(x)\cos'(x)\bigr)/ \cos^2(x) \cr
&= \bigl(\cos^2(x) + \sin^2(x)\bigr) / \cos^2(x) \cr
&= 1 / \cos^2(x) \cr
&= \sec^2(x)
}
$$
This is not valid for $x= \frac{\pi}{2}$; in fact, it is not valid for any $x$ of the form $x = k\,\frac{\pi}{2}$ ($k$ an odd integer). This is because
$\tan$ is not defined at these values.
  
Just as was the case with the logarithmic functions, there are inverses of each of the 
trig functions. Some care must be taken to describe the domains of such 
inverses, but the approach to compute the derivatives of these inverses is the 
same as with the logarithms. We ask the reader to consult a calculus tome 
for a discussion and listing of these functions and their derivatives.


\section{An Interpretation of the Derivative}
Let $d$ and $v$ be the functions which give the distance and 
velocity, respectively, of a car as a function of time. Let's start by 
assuming that the car is traveling at a constant $30$ mph. We know from simple 
physics that when the velocity is constant the 
distance function is 
given as $d(t) = v_0 t$, where $v_0$ is the velocity. This is simply the 
statement distance = rate $\times$ time. In our case this becomes:
$d(t) = 30 t$.
What do the graphs of $d$ and $v$ look like?
\epsfigure{eps/constdist.eps}
{Distance graph when car is traveling $30$ mph.}
\epsfigure{eps/constvelc.eps}
{Velocity graph when car is traveling $30$ mph.}
Another way to say this is that the velocity = distance $\div$ time.
This last statement translates geometrically as follows: The slope of the 
distance graph is $30$. More generally, for a car traveling at a constant 
velocity $v_0$, the slope of the distance graph is $v_0$. That is, at any 
point on the distance graph, $(t, d(t))$, its slope is the velocity at 
time $t$, $v(t)$.

What about more general velocity and distance functions. Can we relate the
distance function, $d(t)$, to the velocity function, $v(t)$ as we did above?
The answer is yes. Take a look at the distance and velocity graphs below. 
\epsfigure{eps/distancepoint.eps}{Distance function.}
\epsfigure{eps/velocitypoint.eps}{Velocity function.}
The \index{tangent line approximation}s to the slope at a point in the distance 
graph are changes in the distance function divided by the time interval. 
This gives the average velocity over that time interval. As the approximations
get better, the intervals of time shrink. In the limit we get the 
velocity at that instant of time; that is
$$
d'(t) = v(t)
$$

Using the same arguments one can show the relation between velocity, $v$,
and \index{acceleration}, $a$:
$$
v'(t) = a(t)
$$

\section{An Application To Physics}
If you throw a ball up in the air with an initial velocity $v_0$ (we take positive values to mean up), how high up
will the ball travel? This is a maximization problem again; we're maximizing the height a ball travels.
From physics, discounting air resistance, the
distance as a function of time is: $d(t) = v_0 t - \frac{1}{2} g t^2$,
where $d(t)$ is the distance in feet at time $t$ in seconds and $g$ is the constant $32$ ft/${\rm sec}^2$.
To solve this problem we need to 
produce \index{candidate input}s for the maximum. This is done by finding the solution 
set, 
$\{ \, t \mid d'(t) = 0 \, \}$, and adding to this set the end points. 
There is only one end point, $0$; we are not constrained by any end time 
condition. The solutions set $\{ \, t \mid d'(t) = 0 \, \}$ is the set of 
solutions to
$d'(t) = (v_0 t)' - (\frac{1}{2} g t^2)' = v_0 - g t = 0$.
There is only one solution which is, $t = v_0/g$. So, the set of input 
candidates is 
$\{\, 0, v_0/g \,\}$. Plugging $0$ in for $t$ gives $d(0) = 0$ as it should.
Plugging in $v_0/g$ in for $t$ gives $d(v_0/g) = v_0^2/g - g v_0^2/(2g^2) =
v_0^2/g - v_0^2/(2g) = v_0^2/(2g)$. If $v_0$ is positive, then this time
corresponds to the maximum height. Positive means throwing the ball up with 
a nonzero speed. So, if we throw the ball up with speed $80$ft/sec, 
(about $55$ mph) then 
the time which maximizes the height is 
$80/g$; this is $80/32 = 2.5$ seconds. Plugging in this time into the 
distance function 
gives a maximum height of $80^2/(2 * 32) = 6400/64 = 100$ feet. 

Where did this distance function come from? Can we derive this formula 
ourselves with what we know? The answer is yes. The earth's gravitation pull 
is constant. This is what Galileo observed: acceleration is constant for 
any object. For our ball this means
$$
a(t) = -g 
$$
To find the corresponding velocity we need to find a function, $v$, such 
that $v' = a$. Notice this is the opposite of differentiation. We need 
to find an {\it anti-derivative\/} of $a$.%
\footnote{\kern 0.5pt \raise 0.5ex \hbox{\dag}}
{This is analogous to finding an {\it \index{anti-difference}\/} of a sequence.}
This is the {\it isolate and invert\/} move of the preface once more:
differentiation is the operation, and anti-differentiation is its inverse.
A function which when differentiated is a constant, $c$, is $c\, t$. 
Therefore, $v(t) = -g\, t$ works. We're using the power rule in reverse. 
One can easily check its derivative 
is $-g$. But is this the only function which when differentiated gives $-g$?
No, \index{anti-derivative}s are not unique. If you remember in an earlier section 
we argued that if $h$ is a function with derivative, $h'$, then 
$h_1$, defined by $h_1(x) = h(x) + c$ also has derivative $h'$. The idea was 
that the constant just shifted the graph of $h$ up or down and didn't change 
$h$'s shape and therefore the slopes of the tangent lines were unchanged.
So, we claim that $v(t) = v_0 - g t$ is also a solution to $v' = a$. In this case, $v(t)$ 
not only has the correct acceleration but also the correct initial velocity.
Now, we also know that $d' = v$. So by the same arguments we need to find 
a function whose derivative is $v_0 - g t$. Using the power rule in reverse 
again, we see that $d$ could be any one of the functions in the family $c + v_0 t - \frac{1}{2} g t^2$. 
The constant $c$ is chosen so that the initial distance is right. Since our distance function is measured
from the ground, initially our distance must be $0$. That is, $d(0) = c + v_0 * 0 - \frac{1}{2} g * 0^2 = c = 0$.
So, $d(t) = v_0 t - \frac{1}{2} g t^2$. We've just derived the distance function used in physics to describe
an object under the influence of gravity -- not a small feat. 


\section{Summary}
We motivated the definition of the derivative through a maximization problem over an interval.
\beginEnum
\enum{Finding the maximum potentially means sifting through infinitely many points -- {\bf this is hard/impossible\/}.}
\enum{Found a {\it geometric\/} way to reduce this search to just a handful of points for smooth functions.
The \index{candidate solution}s come from the end points of the interval and the points where the 
graph of the function is flat.}
\enum{Defined the derivative as a measure of steepness of the graph - this measure is 0 at the
flat points.}
\enum{Found an {\it algebraic\/} way to compute the \index{derivative} (at {\it all\/} points) for a certain class of functions.}
\enum{{\it Extended\/} this class to a much larger one by finding formulas for the derivative of 
linear combinations, products, division, or composition of functions we already know how to differentiate.}
\enum{Now finding the maximum of a {\it nice\/} function can be found from a handful of candidate values:
the end points and the points where $f'(x) = 0$ -- {\bf this is often easy\/}.}
\endEnum

Finally, notice that this whole chapter has been the {\it same game, bigger
board\/} theme of the preface at work -- each tool of the Discrete Calculus
has a continuous counterpart:

\bigskip
\centerline{\vbox{\halign{\quad\hfil#\hfil\quad&\quad\hfil#\hfil\quad\cr
{\bf Discrete}&{\bf Continuous}\cr
\noalign{\smallskip\hrule\smallskip}
$\Delta x_n = x_{n+1} - x_n$&$f'(x) = \lim\limits_{n \rightarrow \infty} \bigl(f(x+h_n) - f(x)\bigr)/h_n$\cr
$\Delta n^k$: leading term $k\,n^{k-1}$&$(x^k)^\prime = k\,x^{k-1}$\cr
linearity of $\Delta$&linearity of differentiation\cr
Discrete Product Rule&Product Rule\cr
anti-difference (summation)&anti-derivative (integration -- next chapter)\cr
}}}
\bigskip


%\section{Other Notions of Derivative}

